%!TEX TS-program = lualatex
%!TEX encoding = UTF-8 Unicode
%! program = lualatex

%%%%% DOCUMENT SETUP %%%%%
% Document class derives from "memoir" class
%\documentclass[10pt, letterpaper, final, twoside, onecolumn, article]{memoir}% For most reports; chapters look like section headers
\documentclass[10pt, letterpaper, final, twoside, onecolumn]{memoir}% For book-like reports with separate chapters

% Useful packages
\usepackage{psrc}
\usepackage{lipsum}
\usepackage{soul}
\usepackage{datetime}
%\usepackage{pgfplots}
\usepackage{pdflscape} % for landscape pages
\usepackage{longtable}
%\usepackage[notes,strict,backend=bibtex8,babel=other,bibencoding=inputenc,doi=false,isbn=false,url=false]{biblatex-chicago}% Chicago bibliography


% Define our harvey balls
\definecolor{cff0000}{RGB}{255,0,0}
\definecolor{cffffff}{RGB}{255,255,255}


\newcommand*\HVHi{
\begin{tikzpicture}[y=0.80pt, x=0.8pt,yscale=-1, inner sep=0pt, outer sep=0pt]
\path[cm={{0.13199701,0.0,0.0,0.13199701,(607.87024,656.47962)}},draw=black,fill=black,miter
  limit=4.00,line width=0.606pt] (700.0000,100.0000) circle (1.4111cm);
\end{tikzpicture}
}

\newcommand*\HHi{
\begin{tikzpicture}[y=0.80pt, x=0.8pt,yscale=-1, inner sep=0pt, outer sep=0pt]
\path[cm={{0.13199701,0.0,0.0,0.13199701,(607.87024,656.47962)}},draw=black,fill=black,miter
  limit=4.00,line width=0.606pt] (550.0000,100.0000) circle (1.4111cm);
\path[fill=cffffff,miter limit=4.00,line width=0.080pt] (680.4686,669.6793) --
  (673.8688,669.6793)arc(180.000:270.000:6.600) -- cycle;
\path[fill=black] (673.8687,669.6793)arc(179.978:0.022:6.600) -- cycle;
\end{tikzpicture}
}

\newcommand*\HMed{
\begin{tikzpicture}[y=0.80pt, x=0.8pt,yscale=-1, inner sep=0pt, outer sep=0pt]
\path[cm={{0.13199701,0.0,0.0,0.13199701,(607.87024,656.47962)}},draw=black,fill=black,miter
  limit=4.00,line width=0.606pt] (400.0000,100.0000) circle (1.4111cm);
\path[fill=cffffff,miter limit=4.00,line width=0.080pt] (660.6690,669.6793) --
  (654.0692,669.6793)arc(180.000:270.000:6.600) -- cycle;
\path[fill=cffffff,miter limit=4.00,line width=0.080pt] (660.6690,669.6793) --
  (654.0692,669.6793)arc(180.000:90.000:6.600) -- cycle;
\end{tikzpicture}
}


\newcommand*\HLow{
\begin{tikzpicture}[y=0.80pt, x=0.8pt,yscale=-1, inner sep=0pt, outer sep=0pt]
\path[cm={{0.13199701,0.0,0.0,0.13199701,(607.87024,656.47962)}},draw=black,fill=cffffff,miter
  limit=4.00,line width=0.606pt] (250.0000,100.0000) circle (1.4111cm);
\path[fill=black,miter limit=4.00,line width=0.000pt] (640.8695,669.6793) --
  (640.8695,663.0795)arc(270.000:360.000:6.600) -- cycle;
\end{tikzpicture}
}

\newcommand*\HVLow{
\begin{tikzpicture}[y=0.80pt, x=0.8pt,yscale=-1, inner sep=0pt, outer sep=0pt]
\path[cm={{0.13199701,0.0,0.0,0.13199701,(607.87024,656.47962)}},draw=black,fill=cffffff,miter
  limit=4.00,line width=0.606pt] (100.0000,100.0000) circle (1.4111cm);
\end{tikzpicture}
}

\title{Scorecard-Based Prioritization of Projects in \\Transportation 2040: Report}
\author{}
%\date{September, 2012}

%\bibliography{PSRC} %Wouldn't a curated, agency-wide bibliography database be swell?

\begin{document}
\mainmatter
\maketitle

\tableofcontents % Want a table of contents?
\clearpage
\listoffigures % I heard you liked these.
\listoftables % Can't have one without the other!
%%%%% BEGIN DOCUMENT %%%%%




%%%% CHAPTER
\chapter{Transportation 2040 Prioritization Background \& Orientation}

\section{Introduction}
This is a working draft document that provides a summary of the proposed “scorecard approach” for evaluating projects in Transportation 2040. This evaluation approach is based on the direction of the PSRC Executive Board on June 28, 2012. Distribution of this document will start with the Regional Staff Committee in September and October 2012, and the Prioritization Working Group in October and November 2012. Refinements will be made to this document based on input from the Regional Staff Committee and the Prioritization Working Group. From there, the draft document will be shared with the PSRC Boards in December 2012. Draft evaluation results begin in Chapter \ref{ch:analysis} of this report.

\section{What is Prioritization?}
Since January 2011, the Transportation Policy Board Working Group, the Region Staff Committee, and PSRC boards and committees have been developing a new framework for evaluating transportation investments in the long-range transportation plan, known as Prioritization.

The Mission Statement of the Transportation 2040 Prioritization Process is:

\begin{quotation}
To prioritize transportation projects in a manner that reinforces and implements VISION 2040. This will include the development and application of an evaluation process to prioritize projects/programs that are within the plan as well as the means for allowing projects to be entered, modified, or removed. (source)
\end{quotation}

This evaluation approach is intended to provide the policy boards with a greater level of detail about the benefits and costs associated with transportation projects. The approach:

\begin{itemize}
\item Addresses State of Good Repair by estimating costs to preserve, maintain, and operate the system by setting a desired regional condition level (e.g. low, medium, or high) for each project category;
\item Excludes significantly funded and voter-approved projects;
\item Separates System Improvements projects into four categories and evaluates these projects in each category by decade;
\item Identifies nine measures (based on VISION 2040 and Transportation 2040) to evaluate System Improvements projects;
\item Uses a scorecard approach with results reported in a simple, graphic format (project cost data will also be reviewed);
\item Reviews projects for consistency with regional policy in VISION 2040 and Transportation 2040.
\end{itemize}

\section{What Must All Projects Do?}
Regardless of a project's mode or category, all projects must:
\begin{itemize}
\item Implement VISION 2040 through a consistent set of measures; and
\item Be consistent with the Regional Growth Strategy.
\end{itemize}

Transportation 2040 implements VISION 2040, the region’s long-range strategy for maintaining a healthy region and promoting the well-being of people and communities, economic vitality, and a healthy environment. VISION 2040 contains 174 multi-county planning policies, 68 actions, and 20 measures to guide long-term development across the region. Transportation 2040 Prioritization depends on the guidance of VISION 2040 and a great deal of effort has gone into making sure that projects are evaluated with a clear reflection of VISION 2040 principles. The graphic below indicates how VISION 2040 will be used in the decision-making process:

The Regional Growth Strategy applies to \emph{all} areas of the four-county region and provides numeric allocations that represent the desired future development pattern for accommodating five million people and 3.1 million jobs by 2040.

Prioritization presents an opportunity to implement local plans and targets with transportation investments. Newly adopted local targets reflect VISION 2040’s Regional Growth Strategy with full implementation occurring through the 2015-2016 comprehensive plan updates.

\section{What Projects will be Prioritized?}

Transportation 2040 Prioritization will apply to projects that are on the Metropolitan Transportation System (MTS) in order to reflect the higher role of regionally significant projects and to maintain local autonomy in decision-making for local projects. Local, non-MTS projects will be sufficiently addressed through a jurisdiction’s comprehensive planning process and will not be prioritized at the regional level. Projects that are within the constrained or unprogrammed parts of Transportation 2040 will also be considered. 

\section{What Projects won't be prioritized?}

Based on direction from the Transportation Policy Board, a certain number of projects in Transportation 2040 will not be a part of Transportation 2040 Prioritization. These projects fall into two categories: significantly funded and voter approved projects.

\subsection{Significantly Funded Projects}

Projects currently in Transportation 2040 that are already sufficiently funded or far enough advanced in the construction process will not be evaluated. Considering these projects in Transportation 2040 Prioritization would not provide useful information for decision-making and, in any event, the outcome of an evaluation would not be able to change the course of project completion. Regional guidance would no longer be useful.

The Prioritization Working Group and Regional Staff Committee created an approach for identifying projects to be excluded from Transportation 2040 Prioritization. The list of projects will be reviewed, and some projects, despite meeting a threshold for exclusion may still be included in the evaluation. The intent is not to simply exclude projects, but rather to consider only those projects for which a review could have a meaningful impact.

The Prioritization Working Group and Regional Staff Committee reached consensus on three tiers of funded projects being removed from further consideration (Table \ref{tb:fundtier}).\footnote{For a complete list of these projects, see Appendix \ref{ch:sigfund}.} The projects in these tiers represent a snapshot in time. As Transportation 2040 Prioritization moves toward the 2014 update to Transportation 2040, additional projects may be added to this list. 
 
\begin{table}[htdp]
\caption{Funding Tiers}
\begin{center}
\begin{tabularx}{0.6\textwidth}{>{\raggedright\arraybackslash}p{3pc}>{\raggedright\arraybackslash}X}
\toprule
TIER 0	&	Approved or conditionally approved projects that are completed \\
TIER 1	&	Approved or conditionally approved projects under construction or with construction soon to begin \\
TIER 2	&	Approved or conditionally approved projects with right of way in progress and construction to start in 1-2 years \\
\bottomrule
\end{tabularx}
\end{center}
\label{tb:fundtier}
\end{table}%
 
 
\subsection{Voter Approved Projects}

The Prioritization Working Group and Regional Staff Committee also agreed on removing voter approved projects from the Transportation 2040 Prioritization process. These projects include Sound Transit 2, Nickel, and Transportation Partnership Program TPA projects. A complete list of these projects is available in Appendix \ref{ch:voterapproved}.


\section{How Are Projects Organized?}
Not all transportation projects are the same, or serve the same purpose. Not every dollar spent yields the same benefits. The requirements for each project may be different based on the project type. To avoid setting unrealistic requirements, projects will be considered within project types or investment categories. For example, a bicycle trail extension project should not be expected to produce large reductions in vehicle travel (as a large transit project might), but bicycle projects should be compared with other bicycle projects in terms of the benefits they provide. Table \ref{tb:projcats} shows the proposed project categories for the Transportation 2040 update to be adopted in 2014.

\begin{table}[htdp]
\caption{default}
\begin{center}
\begin{tabularx}{0.75\textwidth}{>{\raggedright\arraybackslash}p{9pc}>{\raggedright\arraybackslash}X}
\toprule
System Improvements	&	Arterials and State Routes; \\
&	Bicycle, Pedestrian, and Complete Streets; \\
&	Highways; \\
&	Transit Expansion \\
\\
State of Good Repair	&	Maintenance and Preservation; \\
&	Current Operations \\
\\
Programmatic Elements	&	Transportation Demand Management, Transportation System Management \& Operations, Safety, and Special Needs \\
\bottomrule
\end{tabularx}
\end{center}
\label{tb:projcats}
\end{table}%

\subsection{System Improvements}
System Improvements refers to transportation investments that expand transportation infrastructure in one of the following categories: Arterials and State Routes; Bicycle, Pedestrian, and Complete Streets; Highways; and Transit Expansion.

\begin{description}
\item[Arterials and State Routes]  Arterials and portions of state routes that function on arterials are considered in this category. Transit investments that are not related to operations and are not major transit capacity enhancements will also be considered in this category.
\item[Bicycle, Pedestrian, and Complete Streets] Projects within this category focus specifically on projects directly related to any of these three improvements. Projects that are broader in scope or that include multiple modes would be considered under other categories (e.g. an arterial expansion that includes bicycle facilities would be considered in the Arterials and State Routes category and may be ranked higher for including non-roadway expansion elements).
\item[Highway] This category of system investment includes improvements to state and interstate highways.
\item[Transit Expansion] This category includes capital investments in high capacity transit such as light rail, commuter rail, stations, park and ride facilities, trains, buses, and ferries.
\end{description}


\subsection{State of Good Repair}

Transportation 2040’s highest priority is to maintain the existing transportation system in a safe and functional manner. Each element of the transportation network discussed in Transportation 2040 has an associated cost for maintenance and operations through the horizon year of the plan. Given that there is little, if any, project-specific information for maintaining these assets in a state of good repair, costs are estimates through various programmatic models. The priority placed on maintaining and operating existing assets and services prior to expansion will be reflected differently from the individual project evaluations described later in this document. A discussion of how the maintenance and preservation estimates have been conducted to date can be found in Appendix \ref{ch:maintpres}.

PSRC policy board members have communicated the need for additional information on the condition of the transportation network when making decisions about future system improvements. Responding to this direction, this approach suggests that PSRC provide additional information on each component of the transportation network, and to the extent possible, incorporate local asset management planning more directly into the regional planning and decision-making process.


\begin{table}[htdp]
\caption{State of Good Repair Categories}
\begin{center}
\begin{tabularx}{0.9\textwidth}{>{\raggedright\arraybackslash}p{11pc}>{\raggedright\arraybackslash}X}
\toprule

Maintenance \& Preservation	&	Roadways; \\
&	Bridges; \\
&	Bicycle and Pedestrian Facilities; \\
&	Transit Facilities and Rolling Stock; \\
&	Ferry Facilities and Fleet \\
\\
Current Operations	&	System Operations (TSM\&O and Safety); \\
&	Transit Services; \\
&	Ferry Services; \\
&	Transportation Demand Management; \\
&	Other (Administration, Debt Service, Drainage, Lighting, etc.) \\

\bottomrule
\end{tabularx}
\end{center}
\label{default}
\end{table}%

\section{Programmatic Elements}

A large portion of Transportation 2040 is currently handled as a programmatic estimate (i.e. Transportation Demand Management, Transportation System Management and Operations, Special Needs, and Safety). Projects in this category will be evaluated in more detail as additional project information becomes available from project sponsors.

The Programmatic Elements include a second grouping of projects related to local investments not currently on the Metropolitan Transportation System. Examples of these projects include local arterials, collectors, neighborhood streets and bridges, and local bus service (Sound Transit bus service, bus rapid transit, and light rail would be included in the System Improvement category and subject to Transportation 2040 Prioritization). These local, programmatic investments would not be subject to Transportation 2040 Prioritization because they emerged from the local comprehensive planning process.

\section{How Are Projects Measured?}

Transportation 2040 Prioritization will evaluate System Improvements projects according to the following nine consistent measures: Air Quality, Freight, Jobs, Multi-Modal, Puget Sound Land and Water, Safety and System Security, Social Equity and Access to Opportunity, Support for Centers, and Travel.

Table \ref{tb:ninemeasures} reviews the purpose for each measure. The questions and points associated with each measure can be reviewed in Appendix \ref{ch:ninemeasures}. Questions were derived directly from VISION 2040 and are meant to encapsulate the 37 policy concepts highlighted in Using VISION 2040 in Decision-Making.

\begin{table}[htdp]
\caption{Nine Scorecard Measures and Their Purpose}
\begin{center}
\begin{tabularx}{\textwidth}{>{\raggedright\arraybackslash}p{8pc}>{\raggedright\arraybackslash}X}
\toprule
Measure	&	Purpose \\
\midrule
Air Quality	&	Addresses reductions in activities that would lead to reductions in emissions \\ \\
Freight	&	Addresses access to freight-related areas, key freight facilities, and system performance related to freight \\ \\
Jobs		&	Addresses improved access to locations with existing and future concentrations of employment, cluster employment, family-wage jobs, and economic foundations \\ \\
Multi-Modal	&	Addresses improvements to nonmotorized and transit facilities as a means of providing alternatives to single-occupancy vehicles; intended to capture the benefits of providing mobility choices \\ \\
Puget Sound Land and Water	&	Addresses environmental protection through consideration of impacts to resource lands and critical areas as well as designs that improve or reduce stormwater runoff \\ \\
Safety and System Security	&	Addresses improvements to facilities where previous design-related accidents have occurred, or facilities integral to system security \\ \\
Social Equity and Access to Opportunity	&	Addresses improved accessibility and reduced negative impacts for specific segments of the population; includes improvements in health as well as connecting areas of low opportunity with areas of high opportunity using the opportunity maps generated by the Growing Transit Communities projects \\ \\
Support for Centers	&	Addresses improved access within centers, including those that are regional designated, as well as locations that contain transit-supportive land uses and are functionally focal points for the communities where they are located \\ \\
Travel	&	Addresses congestion and efficiency by identifying projects that improve identified existing or future facilities with congestion issues, as well as tools that improve efficiency and throughput \\ \\

\bottomrule
\end{tabularx}
\end{center}
\label{tb:ninemeasures}
\end{table}%


\section{Scorecard Evaluation Approach}

The results of project scoring will demonstrate the relative rank of projects within their respective category as well as how the nine individual measures affect their overall ranking, with the results combined into a total VISION 2040 benefit score. In addition to the scorecard evaluation approach, the travel demand model will be used to provide further information regarding roughly two dozen of the largest highway-related investments in Transportation 2040. The travel demand model results are intended to augment the scorecard evaluation by providing information that PSRC’s analytical tools are particularly well-suited to analyze and not otherwise available through the scorecard evaluation approach. Figure \ref{fg:scexample} shows an example of a potential graphic representation.

\begin{figure}[htbp]
\begin{center}
\includegraphics[width=0.9\textwidth]{figures/scorecard_example}
\caption{Example Graphic Representation of Scorecard Results}
\label{fg:scexample}
\end{center}
\end{figure}

In addition to the scorecard example, results of Transportation 2040 Prioritization will be displayed on a chart that plot the numerical store of each prioritized project along the y-axis against the number of measures that the project scored points in along the x-axis. Project costs will be reflected by the size of the bubble representing the projects on the grid.

[PLACEHOLDER FOR BUBBLE CHART GRAPH]

\section{Cross-Checks}
Once System Improvements projects are evaluated and scored, they will be submitted to three cross-checks---a Regional Growth Strategy cross-check, a Financial cross-check, and a Programmatic cross-check---against other key elements of VISION 2040 and Transportation 2040.

\subsection{Regional Growth Strategy Cross-Check}

The Regional Growth Strategy is based on agreed-upon levels of growth in all regional geographies and represents adopted policy that emphasizes growth in cities, urban areas, and regional and local centers. The Regional Growth Strategy cross-check helps ensure that the benefits of prioritized projects are distributed in a manner roughly proportional to the VISION 2040 Regional Growth Strategy. The cross-check is conducted on the list of System Improvements projects, comparing the distribution of investments and their benefits with how well they align with the distribution of population and employment growth detailed in the Regional Growth Strategy.

\subsection{Financial Cross-Check}

Transportation 2040 includes future revenue assumptions based on a combination of current law revenue, historic trends, and assumption of future user fees. The purpose of an updated financial plan is to provide a financial cross-check that compares current conditions to what was planned when Transportation 2040 was adopted in 2010. A description of funding sources and potential eligible investment options will be outlined for each funding source.

\subsection{Programmatic Cross-Check}

A large portion of Transportation 2040 is currently handled as a programmatic investment—Transportation Demand Management, Transportation System Management and Operations, Special Needs Transportation, and Safety. The Programmatic cross-check will provide a comparison of the realities of implementation versus the direction called for in Transportation 2040. A summary of proposed information to be presented through the Programmatic cross-check is summarized in Table \ref{tb:crosschecks}.
\begin{table}[htdp]
\caption{Programmatic Crosscheck Contents}
\begin{center}
\begin{tabularx}{0.8\textwidth}{>{\raggedright\arraybackslash}p{10pc}>{\raggedright\arraybackslash}X}
\toprule
Investment Type	&	Information Provided \\ \\
\midrule
Safety	&	Target Zero information by mode; comparing trends in the region since 2002 \\ \\
Special Needs Transportation	&	Population change of special needs populations (age, disability, or physical condition) \\ \\
Transportation Demand Management	&	Quantitative and qualitative assessment of program advancements and/or reductions in the region \\ \\
Transportation System Management \& Operations	&	Quantitative and qualitative assessment of program advancements and/or reductions in the region \\
\bottomrule
\end{tabularx}
\end{center}
\label{tb:crosschecks}
\end{table}%




%%%% CHAPTER
\chapter{Transportation 2040 Project Analysis}
\label{ch:analysis}
At the time of this report (\today \ at \currenttime), the status of projects in the sponsor form database are as follows.


\begin{table}[htdp]
\caption{Projects by scorecard reporting status}
\begin{center}
\begin{tabularx}{0.35\textwidth}{>{\raggedright\arraybackslash}X>{\raggedright\arraybackslash}X}
\toprule

\input{tables/projects_by_status.tex}

\end{tabularx}
\end{center}
\label{tb:projects_by_status}
\end{table}%


\begin{table}[htdp]
\scriptsize %small, but better than breaking across page boundaries.
\caption{Project status by jurisdiction}
\begin{center}
\begin{tabularx}{0.5\textwidth}{>{\raggedright\arraybackslash}p{8pc}>{\raggedright\arraybackslash}X>{\raggedright\arraybackslash}X>{\raggedright\arraybackslash}X}
\toprule

\input{tables/projects_by_jurisdiction.tex}

\end{tabularx}
\end{center}
\label{tb:projects_by_jurisdiction}
\end{table}%


 
\chapter{Scoring results}
\lipsum[1-2]

\begin{figure}[htbp]
\begin{center}
\includegraphics[width=\textwidth]{figures/distribution_by_score}

\caption{Distribution of projects by score}
\label{fig:distribution_by_score}
\end{center}
\end{figure}



\begin{figure}[htbp]
\begin{center}
\includegraphics[width=\textwidth]{figures/projects_by_cats_and_score}

\caption{A none-too-useful plot showing the near perfect linear relationship between project score and the number of categories in which the project received points}
\label{fig:projects_by_cats_and_score}
\end{center}
\end{figure}



 
\begin{figure}[htbp]
\begin{center}
\includegraphics[width=\textwidth]{figures/projects_by_score_by_scorecard}

\caption{Distribution of projects by score and by scorecard}
\label{fig:projects_by_score_by_scorecard}
\end{center}
\end{figure}


%%%% Enter landscape mode to make room for wide tables
\clearpage
\begin{landscape}
\pagestyle{empty}


%%%% Non-motorized projects
{\scriptsize
\begin{longtable}{>{\raggedright\arraybackslash}p{5.5pc}>{\raggedright\arraybackslash}p{10.5pc}>{\raggedright\arraybackslash}p{2.5pc}>{\raggedright\arraybackslash}p{2.5pc}>{\raggedright\arraybackslash}p{2.5pc}>{\raggedright\arraybackslash}p{2.5pc}>{\raggedright\arraybackslash}p{2.5pc}>{\raggedright\arraybackslash}p{2.5pc}>{\raggedright\arraybackslash}p{2.5pc}>{\raggedright\arraybackslash}p{2.5pc}>{\raggedright\arraybackslash}p{2.5pc}>{\raggedright\arraybackslash}p{2.5pc}}

\caption{ Non-motorized projects} \\

\toprule

Sponsor	&	Description  & Air Quality & Freight & Jobs & Multi-Modal & Puget Sound Land and Water & Safety and System Security & Social Equity and Opportunity & Support for Centers & Travel & Total Score \\ \midrule
\endfirsthead

\toprule
Sponsor	&	Description  & Air Quality & Freight & Jobs & Multi-Modal & Puget Sound Land and Water & Safety and System Security & Social Equity and Opportunity & Support for Centers & Travel & Total Score \\ \midrule

\endhead

\bottomrule
\multicolumn{12}{r}{\HVHi = Largest benefit, \HHi = Larger benefit, \HMed = Average benefit, \HLow = Smaller benefit, \HVLow = Smallest benefit}
\endfoot

\bottomrule
\multicolumn{12}{r}{\HVHi = Largest benefit, \HHi = Larger benefit, \HMed = Average benefit, \HLow = Smaller benefit, \HVLow = Smallest benefit}
\endlastfoot

\input{tables/tallies_z_nonmotor.tex}

\end{longtable}}




%%%%% Transit and ferry projects
{\scriptsize
\begin{longtable}{>{\raggedright\arraybackslash}p{5.5pc}>{\raggedright\arraybackslash}p{10.5pc}>{\raggedright\arraybackslash}p{2.5pc}>{\raggedright\arraybackslash}p{2.5pc}>{\raggedright\arraybackslash}p{2.5pc}>{\raggedright\arraybackslash}p{2.5pc}>{\raggedright\arraybackslash}p{2.5pc}>{\raggedright\arraybackslash}p{2.5pc}>{\raggedright\arraybackslash}p{2.5pc}>{\raggedright\arraybackslash}p{2.5pc}>{\raggedright\arraybackslash}p{2.5pc}>{\raggedright\arraybackslash}p{2.5pc}}

\caption{Transit projects} \\

\toprule

Sponsor	&	Description  & Air Quality & Freight & Jobs & Multi-Modal & Puget Sound Land and Water & Safety and System Security & Social Equity and Opportunity & Support for Centers & Travel & Total Score \\ \midrule
\endfirsthead

\toprule
Sponsor	&	Description  & Air Quality & Freight & Jobs & Multi-Modal & Puget Sound Land and Water & Safety and System Security & Social Equity and Opportunity & Support for Centers & Travel & Total Score \\ \midrule

\endhead

\bottomrule
\multicolumn{12}{r}{\HVHi = Largest benefit, \HHi = Larger benefit, \HMed = Average benefit, \HLow = Smaller benefit, \HVLow = Smallest benefit}
\endfoot

\bottomrule
\multicolumn{12}{r}{\HVHi = Largest benefit, \HHi = Larger benefit, \HMed = Average benefit, \HLow = Smaller benefit, \HVLow = Smallest benefit}
\endlastfoot%
\input{tables/tallies_z_transit.tex}

\end{longtable}}

%%%%% Arterial projects
{\scriptsize
\begin{longtable}{>{\raggedright\arraybackslash}p{5.5pc}>{\raggedright\arraybackslash}p{10.5pc}>{\raggedright\arraybackslash}p{2.5pc}>{\raggedright\arraybackslash}p{2.5pc}>{\raggedright\arraybackslash}p{2.5pc}>{\raggedright\arraybackslash}p{2.5pc}>{\raggedright\arraybackslash}p{2.5pc}>{\raggedright\arraybackslash}p{2.5pc}>{\raggedright\arraybackslash}p{2.5pc}>{\raggedright\arraybackslash}p{2.5pc}>{\raggedright\arraybackslash}p{2.5pc}>{\raggedright\arraybackslash}p{2.5pc}}

\caption{Arterial roadway projects} \\

\toprule

Sponsor	&	Description  & Air Quality & Freight & Jobs & Multi-Modal & Puget Sound Land and Water & Safety and System Security & Social Equity and Opportunity & Support for Centers & Travel & Total Score \\ \midrule
\endfirsthead

\toprule
Sponsor	&	Description  & Air Quality & Freight & Jobs & Multi-Modal & Puget Sound Land and Water & Safety and System Security & Social Equity and Opportunity & Support for Centers & Travel & Total Score \\ \midrule

\endhead

\bottomrule
\multicolumn{12}{r}{\HVHi = Largest benefit, \HHi = Larger benefit, \HMed = Average benefit, \HLow = Smaller benefit, \HVLow = Smallest benefit}
\endfoot

\bottomrule
\multicolumn{12}{r}{\HVHi = Largest benefit, \HHi = Larger benefit, \HMed = Average benefit, \HLow = Smaller benefit, \HVLow = Smallest benefit}
\endlastfoot

\input{tables/tallies_z_arterial.tex}

\end{longtable}}


%%%%% Projects on state routes
\clearpage
{\scriptsize
\begin{longtable}{>{\raggedright\arraybackslash}p{5.5pc}>{\raggedright\arraybackslash}p{10.5pc}>{\raggedright\arraybackslash}p{2.5pc}>{\raggedright\arraybackslash}p{2.5pc}>{\raggedright\arraybackslash}p{2.5pc}>{\raggedright\arraybackslash}p{2.5pc}>{\raggedright\arraybackslash}p{2.5pc}>{\raggedright\arraybackslash}p{2.5pc}>{\raggedright\arraybackslash}p{2.5pc}>{\raggedright\arraybackslash}p{2.5pc}>{\raggedright\arraybackslash}p{2.5pc}>{\raggedright\arraybackslash}p{2.5pc}}

\caption{Projects on State Routes} \\

\toprule

Sponsor	&	Description  & Air Quality & Freight & Jobs & Multi-Modal & Puget Sound Land and Water & Safety and System Security & Social Equity and Opportunity & Support for Centers & Travel & Total Score \\ \midrule
\endfirsthead

\toprule
Sponsor	&	Description  & Air Quality & Freight & Jobs & Multi-Modal & Puget Sound Land and Water & Safety and System Security & Social Equity and Opportunity & Support for Centers & Travel & Total Score \\ \midrule

\endhead

\bottomrule
\multicolumn{12}{r}{\HVHi = Largest benefit, \HHi = Larger benefit, \HMed = Average benefit, \HLow = Smaller benefit, \HVLow = Smallest benefit}
\endfoot

\bottomrule
\multicolumn{12}{r}{\HVHi = Largest benefit, \HHi = Larger benefit, \HMed = Average benefit, \HLow = Smaller benefit, \HVLow = Smallest benefit}
\endlastfoot

\input{tables/tallies_z_stateroute.tex}

\end{longtable}}


%
%%%%% Projects of unknown category
%%\clearpage  %There are so few, let's not waste another page.
%{\scriptsize
%\begin{longtable}{>{\raggedright\arraybackslash}p{3.5pc}>{\raggedright\arraybackslash}p{3.5pc}>{\raggedright\arraybackslash}p{3.5pc}>{\raggedright\arraybackslash}p{3.5pc}>{\raggedright\arraybackslash}p{3.5pc}>{\raggedright\arraybackslash}p{3.5pc}>{\raggedright\arraybackslash}p{3.5pc}>{\raggedright\arraybackslash}p{3.5pc}>{\raggedright\arraybackslash}p{3.5pc}>{\raggedright\arraybackslash}p{3.5pc}>{\raggedright\arraybackslash}p{3.5pc}>{\raggedright\arraybackslash}p{3.5pc}>{\raggedright\arraybackslash}p{3.5pc}>{\raggedright\arraybackslash}p{3.5pc}}
%
%\caption{Projects of unknown category} \\
%
%\toprule
%
%ID  & Description  & Air Quality & Freight & Jobs & Multi-Modal & Puget Sound Land and Water & Safety and System Security & Social Equity and Access to Opportunity & Support for Centers & Travel & TotalScore \\ \midrule
%\endfirsthead
%
%\toprule
%ID  & Description  & Air Quality & Freight & Jobs & Multi-Modal & Puget Sound Land and Water & Safety and System Security & Social Equity and Access to Opportunity & Support for Centers & Travel & TotalScore \\ \midrule
%\endhead
%
%\bottomrule
%\multicolumn{12}{r}{\HVHi = Largest benefit, \HHi = Larger benefit, \HMed = Average benefit, \HLow = Smaller benefit, \HVLow = Smallest benefit}
%
%\endfoot
%
%\bottomrule
%\multicolumn{12}{r}{\HVHi = Largest benefit, \HHi = Larger benefit, \HMed = Average benefit, \HLow = Smaller benefit, \HVLow = Smallest benefit}
%
%\endlastfoot
%
%\input{tables/tallies_z_unknown.tex}
%
%\end{longtable}}


%%%% Exit landscape mode and return pagestyle to "headings"
\end{landscape}
\pagestyle{headings}


\chapter{Observations \& Recommendations}


\chapter{Proposed Next Steps}

\appendix
\chapter{Significantly Funded Projects}
\label{ch:sigfund}


\chapter{Voter Approved Projects}
\label{ch:voterapproved}

\chapter{State of Good Repair}
\label{ch:maintpres}

\chapter{Consistent Measures Scoring}
\label{ch:ninemeasures}

\chapter{Regional Growth Strategy}

\end{document}